% Standalone problem document, generated from the adjacent README.md.
% Compile with XeLaTeX. Mathematical content is maintained in Markdown.
\documentclass[11pt,a4paper]{article}
\usepackage[fontsize=10.5pt]{scrextend}
\usepackage[margin=22mm,headheight=15pt,headsep=9mm,footskip=11mm]{geometry}
\usepackage{fontspec}
\setmainfont{texgyrepagella}[Extension=.otf,UprightFont=*-regular,BoldFont=*-bold,ItalicFont=*-italic,BoldItalicFont=*-bolditalic]
\setsansfont{texgyreheros}[Extension=.otf,UprightFont=*-regular,BoldFont=*-bold,ItalicFont=*-italic,BoldItalicFont=*-bolditalic]
\setmonofont{texgyrecursor}[Extension=.otf,UprightFont=*-regular,BoldFont=*-bold,ItalicFont=*-italic,BoldItalicFont=*-bolditalic,Scale=MatchLowercase]
\usepackage{amsmath,amssymb}
\usepackage{unicode-math}
\setmathfont{texgyrepagella-math.otf}
\usepackage{xcolor,microtype,enumitem,needspace,fancyhdr}
\usepackage{bookmark}
\definecolor{ink}{HTML}{17344B}
\definecolor{accent}{HTML}{197D80}
\definecolor{muted}{HTML}{596773}
\hypersetup{colorlinks=true,linkcolor=accent,urlcolor=accent,pdftitle={PF-04 - The
maximum cp-rank in order
six},pdfauthor={Open Problems in Numerical Linear Algebra}}
\urlstyle{same}
\setlength{\parindent}{0pt}
\setlength{\parskip}{5pt plus 1pt minus 1pt}
\linespread{1.04}
\setlength{\emergencystretch}{3em}
\widowpenalty=10000
\clubpenalty=10000
\displaywidowpenalty=10000
\allowdisplaybreaks[1]
\setlist{nosep,leftmargin=1.5em}
\providecommand{\tightlist}{\setlength{\itemsep}{0pt}\setlength{\parskip}{0pt}}
\providecommand{\pandocbounded}[1]{#1}
\pagestyle{fancy}
\fancyhf{}
\fancyhead[L]{\sffamily\footnotesize\color{muted}OPEN PROBLEMS IN NUMERICAL LINEAR ALGEBRA}
\fancyhead[R]{\sffamily\bfseries\footnotesize\color{accent}PF-04}
\fancyfoot[L]{\parbox[b]{0.9\textwidth}{\sffamily\fontsize{8}{10}\selectfont\color{muted}Editorial ratings; a bounded search cannot certify that a problem remains unsolved.\\Literature check: 2026-09-10}}
\fancyfoot[R]{\sffamily\footnotesize\color{muted}\thepage}
\renewcommand{\headrulewidth}{0.3pt}
\renewcommand{\headrule}{\hbox to\headwidth{\color{accent}\leaders\hrule height \headrulewidth\hfill}}
\makeatletter
\renewcommand{\section}{\Needspace{5\baselineskip}\@startsection{section}{1}{0pt}{12pt plus 2pt minus 2pt}{4pt}{\sffamily\bfseries\large\color{ink}}}
\renewcommand{\subsection}{\Needspace{4\baselineskip}\@startsection{subsection}{2}{0pt}{9pt plus 1pt minus 1pt}{3pt}{\sffamily\bfseries\normalsize\color{ink}}}
\makeatother
\setcounter{secnumdepth}{0}
\begin{document}
{\sffamily\small\color{muted}Nonnegative and positive
factorizations\par}
\vspace{3pt}
{\raggedright\hyphenpenalty=10000\exhyphenpenalty=10000\sffamily\bfseries\fontsize{21}{24}\selectfont\color{ink}The
maximum cp-rank in order six\par}
\vspace{7pt}
{\sffamily\small\textbf{Difficulty:} challenging\quad\textbf{Importance:} interesting
to the community\par}
{\sffamily\small\bfseries\color{accent}Status: Partially resolved\par}
\vspace{4pt}
{\color{accent}\hrule height 0.8pt}
\vspace{6pt}
\textbf{Rating rationale:} Challenging because sharp cp-rank control is
still missing for the remaining order-six boundary configurations;
community importance is the size of exact completely positive
certificates.

\textbf{Area:} size of nonnegative symmetric factorizations

\subsection{Context and notation}\label{context-and-notation}

A symmetric matrix \(A\) is \textbf{completely positive} if
\(A=BB^\mathsf T\) for an entrywise nonnegative real matrix \(B\) with
finitely many columns. Let \(\mathcal{CP}_n\) be the cone of these
matrices of order \(n\). Its \textbf{cp-rank},
\(\operatorname{cpr}(A)\), is the minimum number of columns in such a
factor, with \(\operatorname{cpr}(0)=0\). Boundaries and interiors use
the usual Euclidean topology on the vector space of real symmetric
matrices.

\subsection{Problem statement}\label{problem-statement}

Is every real completely positive matrix of order six a sum of at most
nine nonnegative rank-one matrices? Equivalently, is

\[
\max_{A\in\mathcal{CP}_6}\operatorname{cpr}(A)=9,
\]

or, equivalently, does each \(A\in\mathcal{CP}_6\) admit
\(A=BB^\mathsf T\) with \(B\in\mathbb R_{\ge0}^{6\times9}\), allowing
zero columns?

The lower bound nine is established. Order six is the unresolved case of
the original Drew--Johnson--Loewy bound: order five satisfies it,
whereas counterexamples exist in higher orders. It is included as the
specific remaining case identified in the literature, rather than as one
item in a dimension-by-dimension list.

A known reduction places a matrix attaining the order-six maximum on the
boundary of \(\mathcal{CP}_6\), with full ordinary rank and at least one
zero entry. This reduction does not establish the nine-column bound.

\subsection{References}\label{references}

Abraham Berman, Mirjam Dür, and Naomi Shaked-Monderer,
\href{https://journals.uwyo.edu/index.php/ela/article/download/1477/1477}{\emph{Open
problems in the theory of completely positive and copositive matrices}},
Electronic Journal of Linear Algebra \textbf{29} (2015), 46--58, §4.2,
p.~53. Naomi Shaked-Monderer,
\href{https://arxiv.org/html/1501.02426v3}{\emph{On the DJL conjecture
for order 6}, corrected arXiv v3}, 2017-06-01; journal version in
Operators and Matrices \textbf{11} (2017), 71--88, Theorem 1.1 and
Corollary 3.2.

\subsection{Status check --- 2026-09-10}\label{status-check-2026-09-10}

Rechecked \href{https://arxiv.org/html/1501.02426v3}{Shaked-Monderer's
corrected v3, Theorem 1.1 and Corollary 3.2}, and searched for later
order-six DJL resolutions. The nine-term bound is proved on the boundary
classes orthogonal to exceptional extremal copositive matrices,
including positive nonsingular boundary matrices. Remaining maximizers
can have a zero entry; this reduction is not a proof for all order-six
matrices. \href{https://doi.org/10.1007/s10957-026-02950-2}{Behling et
al.~(2026), §5.6}, reproduces counterexamples in orders seven and
twelve, not six. No full resolution was located, and the exact
open-status source remains historical.
\end{document}
